\chapter{Related Work}
\label{ch:related}

The \textsc{mathlib5} design sits at the intersection of several mature
research traditions. This chapter positions the work against each, noting
where it adopts established practice and where it extends it.

\section{Formal Verification and Proof Assistants}
\label{sec:rw-provers}

Interactive theorem provers --- Coq~\cite{coq}, Isabelle/HOL~\cite{isabelle},
and Lean~\cite{lean4} --- provide foundational guarantees that a formalized
statement holds. The \textsc{mathlib} project~\cite{mathlib} demonstrates the
viability of large-scale formal mathematics in Lean. \textsc{mathlib5} builds
directly on Lean~4 (Layer~4) and reuses its tactic language, but its
contribution is \emph{not} a new library of theorems. Rather, it uses Lean as
one boundary in a pipeline: obligations are generated upstream (Layers~1--3)
and discharged (or logged) here, then handed to an execution layer.

\section{Refinement Types}
\label{sec:rw-refinement}

Refinement types equip a base type with a logical predicate that values must
satisfy. Liquid Haskell~\cite{liquidhaskell} infers and checks refinement
types via SMT, allowing specifications such as ``this list is sorted'' or
``this index is in bounds'' to be enforced statically. \textsc{mathlib5}
adopts Liquid Haskell at Layer~3 to encode structural constraints on symbolic
expressions (e.g., well-formedness of the AST, non-negativity of degrees). The
bridge in \texttt{layers/liquid/src/Bridge.hs} translates an \textsc{apl}
expression into a Lean proof obligation, making the refinement layer the
generator for Layer~4.

\section{Verified Low-Level Computation}
\label{sec:rw-cmm}

CompCert~\cite{compcert} proved that a C compiler could be verified end to
end. C{--}~\cite{cmm} is a portable, typed assembly language used as a target
for verified backends; it deliberately excludes unstructured control flow and
exposes a clean memory model. \textsc{mathlib5} expresses its numeric kernel
in C{--} (\texttt{kernel/verified\_kernel.cm}) so that the execution layer is
a small, auditable surface. Where CompCert verifies the \emph{compiler},
\textsc{mathlib5} verifies the \emph{kernel primitives} and pairs them with a
separate prover via FFI.

\section{Symbolic and Array Notation}
\label{sec:rw-apl}

\textsc{apl}~\cite{apl} and its descendants (J, K, BQN) compress array
mathematics into dense, unambiguous notation. Iverson's thesis --- that
notation is a tool of thought --- underpins Layer~1. The parser in
\texttt{layers/apl/src/Parser.hs} is a Megaparsec grammar over a small
expression language; it is not a full \textsc{apl} implementation but a
deliberately tractable fragment sufficient to seed the pipeline.

\section{Provenance and Tamper-Evidence}
\label{sec:rw-provenance}

Write-once-read-many (WORM) storage, certificate transparency
logs~\cite{ct}, and git's content-addressed object model all establish
tamper-evident histories. Bitcoin~\cite{bitcoin} extends this with
economic finality: an \texttt{OP\_RETURN} output can anchor an arbitrary
hash into the most adversarial ledger on earth. \textsc{mathlib5} adopts the
WORM pattern for its internal receipt ledger (Layer~9) and, for the
manuscript itself, anchors its fingerprint into Bitcoin (see
\cref{ch:security}). This mirrors Sigstore's~\cite{sigstore} use of
transparency logs for software supply-chain signing, substituting Ed25519 for
Sigstore's Fulcio-issued certificates.

\section{P/NP and Solver Swarms}
\label{sec:rw-pnp}

The observation that \emph{finding} a solution is hard while \emph{checking}
one is easy is the foundation of the P/NP question. \textsc{mathlib5}
operationalizes this at Layer~8: agents compete to produce witnesses
(NP-hard search) while the repository only accepts P-time verifiable proofs.
This is a practical instance of a proof-carrying / proof-search split, related
to formalized contest mathematics and to distributed theorem proving.

\section{Synthesis}
\label{sec:rw-synthesis}

No prior system, to our knowledge, composes \textsc{apl} notation, a canonical
IR, Liquid Haskell refinement, Lean~4 obligations, a C{--} verified kernel,
static analysis, answer-set policy reasoning, a P/NP swarm, WORM provenance,
and automated completeness checking into a single inspectable pipeline sealed
by Ed25519 and Bitcoin. That composition --- and the discipline of treating
each boundary as independently auditable --- is the contribution of
\textsc{mathlib5}.

\section{Two Further Systems}
\label{sec:rel-further}

\paragraph{Isabelle/HOL.} Isabelle shares Lean's goal of a large certified
library but uses higher-order logic and a different tactic language; it has
no built-in numeric closed-form generator for the APL fragment we target.

\paragraph{F* and KreMLin} F* verifies programs in a dependently typed
ML and extracts them to C via KreMLin. Like CertiCoq it is a
verified-extraction system but not a verified numeric-kernel with on-chain
provenance, which is MATHLIB5's distinguishing combination.
